\documentclass[10pt,twocolumn,letterpaper]{article}

%\usepackage[review]{cvpr}      % To produce the REVIEW version
%\usepackage[rebuttal]{cvpr}    % To produce a REBUTTAL
\usepackage{cvpr}              % To produce the CAMERA-READY version
\usepackage{graphicx}
\usepackage{amsmath}
\usepackage{amssymb}
\usepackage{booktabs}


\usepackage[pagebackref,breaklinks,colorlinks]{hyperref}


\def\cvprPaperID{*****} % *** Enter the CVPR Paper ID here
\def\confName{CVPR}
\def\confYear{2023}

\begin{document}
\title{Real-Time FPGA Implementation of a QPSK Modulator and Demodulator with Simulated Channel}
\author{Darsh Gajare, Parth Achat, Nishid Khandagre and Aditya Mali\\
Electrical Engineering\\
Indian Institue of Technology Hyderabad}
\maketitle


\begin{abstract}
	Quadrature Phase Shift Keying (QPSK) is a digital modulation technique which transmits two bits per symbol without additional bandwidth. This project implements a complete, real-time QPSK transmitter and receiver entirely on a Spartan-7 FPGA (Boolean Board, XC7S50). The system performs bit-to-symbol mapping, carrier generation, quadrature mixing, and matched low-pass filtering, and includes a simulated AWGN channel with live-adjustable noise level for testing. System performance is validated through bit-error-rate (BER) measurement and live constellation visualization streamed over UART. This report defines the problem statement, reviews relevant literature on QPSK theory and existing FPGA-based QPSK implementations, and outlines the architecture and verification plan for the mid-term and final phases. Carrier recovery via a Costas loop, symbol timing recovery via the Gardner algorithm, and a coded-versus-uncoded BER comparison are being considered as optional extensions if time permits.


\end{abstract}

\section{Introduction}
\label{sec:intro}

Modern communication systems rely heavily on digital modulation to
move information reliably and efficiently over a band-limited
channel. Digital modulation schemes vary properties of a carrier, such
as its amplitude, phase, or frequency, to represent information while
making efficient use of the available communication bandwidth.

Quadrature Phase Shift Keying (QPSK) is one such scheme. Unlike Binary
PSK, which uses two carrier phases to represent one bit per symbol,
QPSK uses four equally spaced carrier phases, so that each transmitted
symbol represents two bits at once \cite{qpsk_nptel}. This is conveniently described
using an in-phase (I) and quadrature (Q) component, which together
represent the two bits of each transmitted symbol. The transmitted
signal can therefore be represented as a combination of cosine and
sine carriers at the same frequency, with the I and Q components
determining the corresponding symbol.

Field-Programmable Gate Arrays (FPGAs) provide reconfigurable digital
hardware rather than a fixed instruction-execution pipeline: operations
such as carrier generation, multiplication, and filtering can each be
realized as dedicated, concurrently operating hardware blocks. This matters for a real-time
communication system, where all stages of modulation and demodulation
must run continuously at the signal's own rate. Prior work has already demonstrated that
a complete QPSK modulator and demodulator can be realized on FPGA
hardware \cite{7777861,6034354}, which
motivates implementing our system on FPGA rather than restricting it
to a Python/MATLAB simulation alone.

Narrowing from this general motivation to our specific system, this
project implements a complete QPSK transmit-receive chain in which a
bit stream is first mapped to QPSK symbols, carriers are generated by
an NCO, the symbols are modulated onto these carriers, and the
resulting signal is passed through a simulated AWGN channel. At the
receiver, the signal is coherently demodulated, passed through an FIR
filter, and a symbol decision block recovers the original bit stream.
Each of these stages is implemented as a distinct hardware block on
the Spartan-7 FPGA (Boolean Board, XC7S50).

To validate this system, we follow a staged verification approach: a
software reference model is developed first, followed by a
fixed-point version of the same algorithm, and finally the FPGA/RTL
implementation itself. At each stage, the implementation's behavior is
compared against the reference model, and the overall system is
evaluated using bit-error-rate (BER), constellation diagrams, behavior
under varying noise conditions, and FPGA implementation
characteristics such as resource utilization.

\section{Literature Review}
\label{sec:litreview}

The most directly relevant prior work is Popescu \etal
\cite{6034354}, who implement a complete QPSK modulator,
channel, and demodulator rather than a modulator alone. Their design
uses two Spartan-3E boards, with one board implementing the modulator
and the other implementing the demodulator. Sine and cosine carriers
are generated internally from a ROM and combined with the odd and even
bit sequences to form the QPSK signal. At the receiver, the received
signal is multiplied by locally regenerated sine and cosine references,
and the recovered signals are compared with a threshold to reconstruct
the transmitted bits. This work demonstrates the feasibility of a
complete FPGA-based QPSK transmit-receive chain and is closely related
to our proposed system. Our implementation differs in that the
complete digital processing chain is targeted to a single Spartan-7
FPGA rather than two separate FPGA boards.

Carrier generation is an important hardware block in a digital PSK
modulator. Al Safi and Bazuin \cite{7777861} present FPGA
implementations of BPSK and QPSK modulators using accumulator-based
addressing and lookup-table techniques. For QPSK, the required
phase-shifted sinusoidal waveforms are generated using accumulators and
a single lookup table, with the implementation carried out in VHDL.
This work is directly relevant to our carrier-generation block, where
a phase-accumulator-driven numerically controlled oscillator (NCO) with
a lookup table is used to generate the sine and cosine references for
modulation and demodulation.

Not every FPGA QPSK design follows the same I/Q-based structure.
Pareek \cite{7509346} compares a conventional I/Q approach with a
digitized-carrier approach. In the proposed digitized approach, the
input bit pair determines the required phase, and a 4:1 multiplexer
selects the corresponding phase of the digitized carrier. Both
architectures are implemented on a Spartan-3E FPGA using Verilog HDL.
The work therefore demonstrates that QPSK modulation can be implemented
using different hardware architectures. Our project adopts the
conventional I/Q and carrier-generation approach rather than the
multiplexer-based phase-selection approach.

Taken together, these works demonstrate FPGA-based QPSK modulation,
complete modulator-channel-demodulator systems, and different
approaches to digital carrier generation and hardware architecture.
The objective of this project is not to introduce a new QPSK algorithm,
but to integrate established techniques into a self-contained and
verifiable FPGA-based system. The proposed implementation targets a
single Spartan-7 XC7S50 FPGA and includes QPSK mapping, NCO-based
carrier generation, modulation, a simulated AWGN channel, coherent
demodulation, FIR filtering, symbol decision, and BER measurement.
Selected I/Q samples will be streamed through UART for constellation
visualization on a host computer.

\section{Next Steps}

{\small
\bibliographystyle{ieee_fullname}
\bibliography{references}
}
\end{document}
